\documentclass[12pt]{article}

\usepackage{ed}
\usepackage{amsmath}

\setImageBase{../classes/myclass/images}

\begin{document}

\begin{topic}{Differentiation conceptual}

  \begin{question}{Multiple Choice}
    \name{Visual recognition 1}

    \qutext{Consider a function $f(x)$ sketched below.

      \graph[12]{(-1)^$n*(x-$a1)*(x-$a2)*(x-$a3)/10}{-6}{6}{-6}{6}

      Which curve best resembles $f'(x)$?}

    \choice*{\graph[12]{(-1)^$n*((x-$a1)*(x-$a2)+(x-$a1)*(x-$a3)+(x-$a2)*(x-$a3))/10}{-6}{6}{-6}{6}}

    \choice{\graph[12]{-(-1)^$n((x-$a1)*(x-$a2)+(x-$a1)*(x-$a3)+(x-$a2)*(x-$a3))/10}{-6}{6}{-6}{6}}

    \choice{\graph[12]{(x-$a1)*(x-$a3)/10}{-6}{6}{-6}{6}}

    \choice{\graph[12]{-(x-$a1)*(x-$a3)/10}{-6}{6}{-6}{6}}

    \code{
      $a1 = range(-5,-2);
      $a2 = range(-1,1);
      $a3 = range(2,5);
      $n = range(1,2);
    }

  \end{question}

\end{topic}

\begin{topic}{Product rule - elementary}

  %% Only basic product rule questions involving polynomials and powers.

  \begin{question}{Formula}
    \qutext{Differentiate $f(x)$ where
      \[
      f(x) = (\var{a1} x^{\var{n1}} + \var{b1}) ( \var{a2} x^{\var{n2}} + \var{b2}).
      \]}

    \answer{\var{tmp1} x^{\var{\var{n1}-1}} ( \var{a2} x^{\var{n2}} +
      \var{b2}) +
      \var{tmp2} x^{\var{\var{n2}-1}}(\var{a1} x^{\var{n1}} + \var{b1})}

    \code{
      $a1 = range(2,6);
      $a2 = range(2,6);
      $n1 = range(3,6);
      $n2 = range(3,6);
      $b1 = range(1,5);
      $b2 = range(1,5);
      $tmp1 = $a1*$n1;
      $tmp2 = $a2*$n2;
    }

  \end{question}

  \begin{question}{Formula}
    \qutext{Differentiate $f(x)$ where
      \[
      f(x) = (\frac{\var{a1}}{x^{\var{n1}}} + \var{b1}) ( \var{a2} x^{\var{n2}} + \var{b2}).
      \]}

    \answer{-\var{tmp1} x^{\var{-\var{n1}-1}} ( \var{a2} x^{\var{n2}} +
      \var{b2}) +
      \var{tmp2} x^{\var{\var{n2}-1}}(\var{a1} x^{\var{-\var{n1}}} + \var{b1})}

    \code{
      $a1 = range(2,6);
      $a2 = range(2,6);
      $n1 = range(3,6);
      $n2 = range(3,6);
      $b1 = range(1,5);
      $b2 = range(1,5);
      $tmp1 = $a1*$n1;
      $tmp2 = $a2*$n2;
    }

  \end{question}

  \begin{question}{Formula}
    \qutext{Differentiate $f(x)$ where
      \[
      f(x) = (\frac{\var{a1}}{x^{\var{n1}}} + \var{b1}) 
      ( \frac{\var{a2}}{x^{\var{n2}}} + \var{b2}).
      \]}

    \answer{-\var{tmp1} x^{\var{-\var{n1}-1}} ( \var{a2} x^{\var{-\var{n2}}} +
      \var{b2}) -
      \var{tmp2} x^{\var{-\var{n2}-1}}(\var{a1} x^{\var{-\var{n1}}} + \var{b1})}

    \code{
      $a1 = range(2,6);
      $a2 = range(2,6);
      $n1 = range(3,6);
      $n2 = range(3,6);
      $b1 = range(1,5);
      $b2 = range(1,5);
      $tmp1 = $a1*$n1;
      $tmp2 = $a2*$n2;
    }

  \end{question}

  \begin{question}{Formula}
    \qutext{Differentiate $f(x)$ where
      \[
      f(x) = (\var{a1} x^{\var{n1}} + \var{b1}) 
      (\sqrt[\var{n2}]{x} + \var{b2}).
      \]}

    \answer{\var{tmp1} x^{\var{\var{n1}-1}} (x^{(\var{n3})} + \var{b2})
      + \var{n3} x^{-(\var{n4})} (\var{a1} x^{\var{n1}} + \var{b1})}

    \code{
      $a1 = range(2,6);
      $a2 = range(2,6);
      $n1 = range(3,6);
      $n2 = range(3,5);
      $b1 = range(1,5);
      $b2 = range(1,5);
      $tmp1 = $a1*$n1;
      $n3 = frac(1,$n2);
      $n4 = frac($n2-1,$n2);
    }

  \end{question}

\end{topic}

\begin{topic}{Quotient rule - elementary}

  %% Only simple quotient rule questions here involving polynomials and
  %% powers. 

  \begin{question}{Formula}
    \qutext{Differentiate $f(x)$ where
      \[
      f(x) = \frac{\var{a1} x + \var{b1}}
      {\var{a2} x + \var{b3}}.
      \]}

    \answer{\var{tmp2} (\var{a2} x + \var{b3})^(-2)}

    %% Watch out for the determinant being zero.
    \code{
      $a1 = range(1,6);
      $a2 = range(1,6);
      $b1 = range(1,6);
      $b2 = range(1,6);
      $tmp1 = int($a1*$b2 - $b1*$a2);
      $b3 = $b2 + eq($tmp1,0);
      $tmp2 = int($b3*$a1 - $b1*$a2);
    }

  \end{question}

  \begin{question}{Formula}
    \qutext{Differentiate $f(x)$ where
      \[
      f(x) = \frac{\var{a1} x^2 + \var{b1}}
      {\var{a2} x + \var{b2}}.
      \]}

    \answer{(\var{\var{a1}*\var{a2}} x^2 + \var{2*\var{a1}*\var{b2}} x -
      \var{\var{b1}*\var{a2}})
      (\var{a2} x + \var{b2})^(-2)}

    \code{
      $a1 = range(1,6);
      $a2 = range(1,6);
      $b1 = range(1,6);
      $b2 = range(1,6);
    }

  \end{question}

  \begin{question}{Formula}
    \qutext{Differentiate $f(x)$ where
      \[
      f(x) = \frac{\var{a1} x + \var{b1}}
      {\var{a2} x^2 + \var{b2}}.
      \]}

    \answer{(-\var{\var{a1}*\var{a2}} x^2 - \var{2*\var{a2}*\var{b1}} x +
      \var{\var{a1}*\var{b2}})
      (\var{a2} x^2 + \var{b2})^(-2)}

    \code{
      $a1 = range(1,6);
      $a2 = range(1,6);
      $b1 = range(1,6);
      $b2 = range(1,6);
    }

  \end{question}

  \begin{question}{Formula}
    \qutext{Differentiate $f(x)$ where
      \[
      f(x) = \frac{\var{a1} x^2 + \var{b1}x}
      {\var{a2} x + \var{b2}}.
      \]}

    \answer{(\var{\var{a1}*\var{a2}} x^2 + \var{2*\var{a1}*\var{b2}} x +
      \var{\var{b1}*\var{b2}})
      (\var{a2} x + \var{b2})^(-2)}

    \code{
      $a1 = range(1,6);
      $a2 = range(1,6);
      $b1 = range(1,6);
      $b2 = range(1,6);
    }

  \end{question}

  \begin{question}{Formula}
    \qutext{Differentiate $f(x)$ where
      \[
      f(x) = \frac{\var{a1} x + \var{b1}}
      {\var{a2} x^2 + \var{b2}x}.
      \]}

    \answer{(-\var{\var{a1}*\var{a2}} x^2 - \var{2*\var{a2}*\var{b1}} x -
      \var{\var{b1}*\var{b2}})
      (\var{a2} x^2 + \var{b2}x)^(-2)}

    \code{
      $a1 = range(1,6);
      $a2 = range(1,6);
      $b1 = range(1,6);
      $b2 = range(1,6);
    }

  \end{question}

\end{topic}

\begin{topic}{Trigonometric functions}

  %% Only basic trig questions here.

  \begin{question}{Formula}

    \qutext{What is the derivative of $\sin x$?}

    \answer{cos x}

  \end{question}

  \begin{question}{Formula}

    \qutext{What is the derivative of $\cos x$?}

    \answer{-sin x}

  \end{question}

  \begin{question}{Formula}

    \qutext{What is the derivative of $\tan x$?}

    \answer{(sec x)^2}

  \end{question}

  \begin{question}{Formula}

    \qutext{What is the derivative of $\cot x$?}

    \answer{-(csc x)^2}

  \end{question}

  \begin{question}{Formula}

    \qutext{What is the derivative of $\sec x$?}

    \answer{(sec x) (tan x)}

  \end{question}

  \begin{question}{Formula}

    \qutext{What is the derivative of $\csc x$?}

    \answer{-(csc x) (cot x)}

  \end{question}

\end{topic}

\begin{topic}{Product rule - trig}

  \begin{question}{Formula}

    \qutext{What is the derivative of $x^{\var{n}} \sin(x)$?}

    \answer{\var{n} x^{\var{nm1}} sin(x) + x^{\var{n}} cos(x)}

    \code{
      $n = range(3,6);
      $nm1 = $n-1;
    }

  \end{question}

  \begin{question}{Formula}

    \qutext{What is the derivative of $(x^{\var{n}}+x) \tan(x)$?}

    \answer{(\var{n} x^{\var{nm1}}+1) tan(x) + (x^{\var{n}}+x) (sec(x))^2}

    \code{
      $n = range(3,6);
      $nm1 = $n-1;
    }

  \end{question}

  \begin{question}{Formula}

    \qutext{What is the derivative of $(x^{\var{n}}+\sqrt{x}) \csc(x)$?}

    \answer{(\var{n} x^{\var{nm1}}+(1/2) x^(-1/2)) csc(x) - 
      (x^{\var{n}}+sqrt(x)) csc(x) cot(x)}

    \code{
      $n = range(3,6);
      $nm1 = $n-1;
    }

  \end{question}

  \begin{question}{Formula}

    \qutext{What is the derivative of $\sin(x) \cos(x)$?}

    \answer{(cos(x))^2 - (sin(x))^2}

  \end{question}

  \begin{question}{Formula}

    \qutext{What is the derivative of $\sin^2(x) \var{pmsw} \cos^2(x)$?}

    \comment{One can use angle sum formulas to simplify this problem
    before differentiating.}

    \answer{$answ}

    \code{
      $sw = range(0,1);
      $pmsw = switch($sw,"+","-");
      $answ = switch($sw,0,"4 sin(x) cos(x)");
    }
  \end{question}

\end{topic}

\begin{topic}{Quotient rule - trig}

  \begin{question}{Formula}

    \qutext{What is the derivative of $\frac{\sin x}{x^{\var{n}}}$?}
    
    \answer{(x^\var{n} cos(x) - \var{n} x^\var{nm1} sin(x))/x^{\var{nt2}}}

    \code{
      $n = range(3,6);
      $nm1 = $n-1;
      $nt2 = 2*$n;
    }
  \end{question}
  
  \begin{question}{Formula}

    \qutext{What is the derivative of $\frac{x^{\var{n}}+x}{\tan x}$?}
    
    \comment{You can use the quotient rule or the product rule here.}
    
    \answer{(\var{n} x^\var{nm1}+1) cot(x) - (x^{\var{n}}+x) (csc(x))^2}

    \code{
      $n = range(3,6);
      $nm1 = $n-1;
      $nt2 = 2*$n;
    }
  \end{question}
  
  \begin{question}{Formula}

    \qutext{What is the derivative of $\frac{\cos x}{x^{\var{n}}+1}$?}
    
    \answer{(-(x^\var{n}+1) sin(x) + \var{n} x^\var{nm1} cos(x))/(x^{\var{n}}+1)^2}

    \code{
      $n = range(3,6);
      $nm1 = $n-1;
      $nt2 = 2*$n;
    }
  \end{question}
  
  \begin{question}{Formula}

    \qutext{What is the derivative of $\frac{\tan x}{\cos (x) +1}$?}
    
    \answer{((sec(x))^2 (cos(x)+1) - tan(x) sin(x))/(cos(x)+1)^2 }
  
  \end{question}
  
  \begin{question}{Formula}

    \qutext{What is the derivative of $\frac{\cos x+1}{\sin x}$?}

    \comment{If you simplify first, you need not use the quotient rule
    at all.}
    
    \answer{-(csc(x))^2 - csc(x) cot(x) }
  
  \end{question}
  
\end{topic}

\end{document}
